\chapter{Testes e valida\c{c}\~ao}
\label{cap:testes}

\section{Estrat\'egia}

O projeto n\~ao possui testes unit\'arios automatizados no sentido usual
(\code{cargo test}), porque o c\'odigo depende de hardware, de um bootloader
e de um ambiente de execu\c{c}\~ao bare metal. A valida\c{c}\~ao foi feita por
\emph{testes de integra\c{c}\~ao} em QEMU, com verifica\c{c}\~ao de:

\begin{itemize}
  \item sa\'ida textual na serial;
  \item conte\'udo da mem\'oria VGA, lido pelo monitor do QEMU;
  \item entrada pelo teclado PS/2, simulada com \code{sendkey};
  \item entrada pela serial, enviada por socket Unix;
  \item validade do ELF e da ISO com ferramentas externas.
\end{itemize}

\section{Matriz de testes}

\begin{table}[H]
\centering
\caption{Matriz de testes executados.}
\label{tab:testmatrix}
\begin{tabularx}{\linewidth}{@{}lXc@{}}
\toprule
\textbf{ID} & \textbf{Teste} & \textbf{Resultado} \\
\midrule
T01 & \code{make clean \&\& make} conclui sem erros & Passou \\
T02 & \code{grub-file --is-x86-multiboot2} aceita o ELF & Passou \\
T03 & ISO reconhecida como El Torito BIOS & Passou \\
T04 & QEMU/SeaBIOS inicializa a ISO e chega ao banner & Passou \\
T05 & Prompt \code{rinux>} aparece na tela VGA & Passou \\
T06 & Entrada PS/2: \code{help} lista os comandos & Passou \\
T07 & Entrada PS/2: \code{echo ola} imprime \code{ola} & Passou \\
T08 & Entrada serial: \code{echo serial ok} funciona & Passou \\
T09 & Prompt retorna ap\'os cada comando & Passou \\
T10 & \code{info} mostra as configura\c{c}\~oes do kernel & Passou \\
T11 & \code{clear} envia \code{ESC[2J ESC[H} e limpa a tela & Passou \\
T12 & \code{reboot} e \code{halt} encerram o fluxo & Verificado manualmente \\
T13 & IDT instalada e timer IRQ0 incrementa \code{ticks} & Passou \\
T14 & \code{demand} gera \code{\#PF} e mapeia a p\'agina do kernel & Passou \\
T15 & \code{usertest} roda em ring 3, gera dois page faults e retorna & Passou \\
T16 & \code{userprot} viola U/S e \'e encerrado pelo tratador & Passou \\
T17 & \code{mem} mostra U/S = 1 para usu\'ario e 0 para o kernel & Passou \\
T18 & \code{tty} reporta \code{tty0} e \code{tty 3} troca de console & Passou \\
T19 & \code{getty} autentica \code{root}/\code{rinux} em ring 3 & Passou \\
T20 & Senha n\~ao \'e ecoada; senha errada encerra com status 1 & Passou \\
\bottomrule
\end{tabularx}
\end{table}

\section{Evid\^encias de execu\c{c}\~ao}

\subsection{Banner e prompt}

\begin{lstlisting}[style=shell,caption={Tela VGA ap\'os o boot.},label={lst:testboot}]
==============================================================
  Rust VGA Kernel - x86_64 / modo longo
  Boot: Multiboot2 (BIOS)
  Video: VGA texto 80x25 @ 0xB8000
  Memoria: paginas de 4 KiB + demand paging
  Interrupcoes: IDT + PIC + timer IRQ0 (~100 Hz)
==============================================================
Multiboot2: magic=0x36d76289 info=0x00112040

Digite 'help' para ver os comandos.

rinux>
\end{lstlisting}

O valor de \code{info} varia a cada execu\c{c}\~ao, porque depende da
aloca\c{c}\~ao feita pelo GRUB. O \code{magic} deve ser sempre
\code{0x36d76289}.

\subsection{Comando \code{help} via PS/2}

\begin{lstlisting}[style=shell,caption={Sa\'ida de \code{help} ap\'os entrada PS/2.},label={lst:testhelp}]
rinux> help
Comandos disponiveis:
  help          mostra esta ajuda
  clear         limpa o console ativo
  echo <texto>  imprime o texto
  info          informacoes do kernel
  tty [0..7]    mostra ou troca o console virtual
  getty [0..7]  roda o login em ring 3 no console
  demand        testa demand paging no kernel
  usertest      roda um programa em ring 3 com page fault
  userprot      testa a protecao usuario/kernel
  mem           estatisticas de frames e paginas
  reboot        reinicia o PC
  halt          para a CPU
rinux>
\end{lstlisting}

\subsection{Comando \code{echo} via PS/2}

\begin{lstlisting}[style=shell,caption={Sa\'ida de \code{echo ola} ap\'os entrada PS/2.},label={lst:testecho}]
rinux> echo ola
ola
rinux>
\end{lstlisting}

\subsection{\code{clear} via ANSI}

\begin{lstlisting}[style=shell,caption={Sequ\^encia de escape observada na COM1.},label={lst:testclear}]
rinux> clear
^[[2J^[[H
rinux>
\end{lstlisting}

A nota\c{c}\~ao com circunflexo \'e como um monitor serial exibe o byte
\code{0x1B}. No console VGA o comando apaga as 2.000 c\'elulas e devolve o
cursor para \code{(0,0)}, sem imprimir os caracteres de escape como texto.

\subsection{Entrada serial}

\begin{lstlisting}[style=shell,caption={Entrada e sa\'ida pela COM1.},label={lst:testserial}]
echo serial ok
serial ok
rinux>
\end{lstlisting}

O texto digitado no socket serial aparece no VGA porque o shell ecoa cada
caractere e o console escreve tanto no VGA quanto na serial.

\section{Script de teste headless}

O procedimento abaixo pode ser adaptado para CI. Ele inicia o QEMU sem
interface gr\'afica, envia um comando pela serial e l\^e a sa\'ida.

\begin{lstlisting}[language=Python,caption={Esqueleto de teste headless com socket serial.},label={lst:testscript}]
import socket, subprocess, time, os

ISO = "vga-kernel.iso"
SER = "/tmp/vga-serial.sock"
MON = "/tmp/vga-mon.sock"

for path in (SER, MON):
    try: os.unlink(path)
    except FileNotFoundError: pass

qemu = subprocess.Popen([
    "qemu-system-x86_64", "-cdrom", ISO, "-boot", "order=d",
    "-m", "128M", "-display", "none",
    "-serial", f"unix:{SER},server=on,wait=off",
    "-monitor", f"unix:{MON},server=on,wait=off",
])

time.sleep(7)
ser = socket.socket(socket.AF_UNIX, socket.SOCK_STREAM)
ser.connect(SER)
ser.sendall(b"echo serial ok\r")
time.sleep(1)
data = ser.recv(65536).decode(errors="replace")
assert "serial ok" in data
ser.close()
qemu.terminate()
qemu.wait(timeout=10)
\end{lstlisting}

Para verificar o VGA, conecte ao monitor, envie
\code{xp /4000bx 0xb8000}, interprete os pares (caractere, atributo) e
reconstrua as linhas. O teste PS/2 usa \code{sendkey}:

\begin{lstlisting}[style=shell,caption={Envio de teclas pelo monitor do QEMU.},label={lst:sendkey}]
sendkey h
sendkey e
sendkey l
sendkey p
sendkey ret
\end{lstlisting}

\begin{atencao}[N\~ao envie \code{quit} antes de terminar]
O comando \code{quit} no monitor encerra o QEMU. Se o script ainda precisar
inspecionar o VGA ou enviar mais teclas, use uma conex\~ao separada para cada
comando e s\'o encerre o processo ao final.
\end{atencao}

\section{Testes de build e formato}

\begin{lstlisting}[style=shell,caption={Comandos de valida\c{c}\~ao estrutural.},label={lst:testbuild}]
make clean && make
grub-file --is-x86-multiboot2 target/x86_64-unknown-none/release/vga-kernel
readelf -h target/x86_64-unknown-none/release/vga-kernel
readelf -l target/x86_64-unknown-none/release/vga-kernel
xorriso -indev vga-kernel.iso -report_el_torito plain
\end{lstlisting}

Resultados esperados:

\begin{itemize}
  \item \code{grub-file} retorna c\'odigo 0;
  \item \code{readelf -h} mostra \code{Type: EXEC}, \code{Machine: Advanced
        Micro Devices X86-64} e \code{Entry point address: 0x100018};
  \item \code{readelf -l} mostra um segmento baixo em \code{0x100000} e
        segmentos altos em \code{0xFFFFFFFF8010...};
  \item \code{xorriso} reporta \emph{El Torito} com
        \code{boot\_hybrid.img}.
\end{itemize}

\section{Limita\c{c}\~oes dos testes}

\begin{itemize}
  \item N\~ao h\'a cobertura automatizada de todos os caminhos do teclado
        (Shift, Tab, Backspace, teclas estendidas).
  \item O teste PS/2 depende do \code{sendkey} do QEMU e do layout de teclas
        do monitor; nem todos os caracteres s\~ao f\'aceis de simular.
  \item A leitura do VGA pelo monitor depende do formato de dump e do
        endere\c{c}o \code{0xB8000}.
  \item O comportamento em hardware f\'isico pode diferir do QEMU,
        especialmente em UEFI sem CSM, teclados USB e chipsets sem 8042.
  \item \code{reboot} e \code{halt} foram verificados manualmente; n\~ao h\'a
        teste automatizado de reset.
  \item N\~ao h\'a testes de estresse, tempo ou concorr\^encia.
\end{itemize}

\section{Checklist de aceita\c{c}\~ao}

\begin{itemize}
  \item[$\square$] \code{make} gera ELF e ISO sem erros.
  \item[$\square$] \code{grub-file} reconhece Multiboot2.
  \item[$\square$] QEMU chega ao banner e ao prompt \code{rinux>}.
  \item[$\square$] \code{help}, \code{echo}, \code{info} e \code{clear}
        funcionam pelo teclado PS/2.
  \item[$\square$] Entrada pela serial funciona.
  \item[$\square$] \code{info} mostra o timer e \code{mem} mostra
        \emph{ticks} crescentes.
  \item[$\square$] \code{demand} mapeia a p\'agina do kernel ap\'os
        \code{\#PF}.
  \item[$\square$] \code{usertest} gera dois page faults e retorna ao shell.
  \item[$\square$] \code{userprot} \'e encerrado pela prote\c{c}\~ao U/S.
  \item[$\square$] \code{tty 3} troca para \code{tty3} e \code{tty 0}
        retorna ao \code{tty0}.
  \item[$\square$] \code{getty} aceita \code{root}/\code{rinux} e volta ao
        shell; senha errada encerra com status 1.
  \item[$\square$] \code{halt} para a CPU e \code{reboot} reinicia (quando o
        ambiente suporta).
  \item[$\square$] A documenta\c{c}\~ao LaTeX compila e gera o PDF.
\end{itemize}
